\plannedchapter{ELF、符号、重定位、链接和加载}
{本章系统讲解 Linux x86-64 二进制文件格式和链接加载流程，为后续 ABI、动态库、profiling 和生产库阅读打基础。}
{\item ELF 文件为什么需要 header、section、segment？
\item symbol、relocation、PLT、GOT 分别解决什么问题？
\item 静态链接和动态链接的性能与工程取舍是什么？}
{\item 目标文件中未解析引用如何通过 relocation 延迟到链接期。
\item 动态链接器如何加载共享库并解析符号。
\item 函数调用经过 PLT/GOT 的路径以及 lazy binding 的基本模型。}
{用 \code{readelf -h/-S/-l/-r}、\code{nm -C}、\code{objdump -drwC} 分析一个含外部函数调用的程序；写静态库和共享库，对比符号和加载行为。}
